% DCA-Trie Journal Paper — Title, Abstract, Keywords
% Session 1: Title + Abstract + Keywords
% Target: 200-250 word abstract, 5-7 keywords

\begin{titlepage}
  \centering
  \vspace*{2cm}

  {\LARGE\bfseries Dynamic Context-Aware Tries for Knowledge Graph-Constrained Large Language Model Generation\par}

  \vspace{1.5cm}

  {
    Bernard Kirk Adjanor Katamanso\textsuperscript{1},
    Erica Amonor\textsuperscript{1},
    Joseph Osei Nyarko\textsuperscript{1},
    Jessica Afua Etornam Nsafuah\textsuperscript{1}

    \vspace{0.5cm}

    \textsuperscript{1}University of Mines and Technology, Tarkwa, Ghana
  }

  \vspace{2cm}

  \begin{abstract}
    Knowledge Graph constrained decoding eliminates hallucination in large language model (LLM) reasoning by restricting output to verified paths. The leading approach, Graph Constrained Reasoning (GCR),achieves zero hallucination on Knowledge Graph Question Answering (KGQA) benchmarks by encoding all structurally valid paths in a static trie. The constraint is absolute but indifferent to the question: a three-hop query with twenty out-degree relations admits roughly 24{,}000 paths, and the LLM must navigate among them using the same internal knowledge that causes hallucination. We introduce  Dynamic Context-Aware Tries (DCA-Trie), a dynamic constrained decoding framework that prunes the valid path set during generation using a TypeOracle, a lightweight symbolic module that checks two ontology gates per candidate triple. On WebQSP, DCA-Trie removes 85.5\% of semantically irrelevant paths while maintaining a false negative rate below 5\%. We further introduce the Semantic Irrelevance Ratio (SIR), a metric that separates constraint quality from answer accuracy, enabling independent evaluation of what the oracle filters versus what the LLM actually uses. SIR reveals a non-monotone relationship between constraint tightness and answer correctness: tighter constraints reduce the search space but do not improve Hits@1. We identify four mechanisms that drive this decoupling and show that the result holds across ablation conditions. DCA-Trie and SIR together provide the first framework for independently evaluating constraint quality and answer accuracy in knowledge graph-constrained decoding.
  \end{abstract}

  \vspace{1cm}

  \noindent\textbf{Keywords:} knowledge graph question answering, constrained decoding, large language models, hallucination mitigation, trie-based decoding, semantic filtering, type oracle

\end{titlepage}
